\section{Linear equations of order $n$ with constant coefficients}

An ordinary differential equation of order $n$ with constant coefficients is an
equation of the form:
\mymark{***}
\begin{equation}\label{eq:edo_lineare_omogenea_coeff_costanti_non_omogenea}
    a_n\cdot u^{(n)}(x) + a_{n-1} \cdot u^{(n-1)}(x) + \dots + a_1 \cdot u'(x) +
  a_0 \cdot u(x) = g(x)
%  \sum_{k=0}^n a_k \cdot u^{(k)}(x) = g(x)
\end{equation}
with $a_0, a_1, \dots a_n \in \RR$ (these are the constant coefficients)
$a_n\neq 0$ and $g(x)$ a continuous function
(if $a_n=0$ it is not a problem: we have an equation of lower order).
Such an equation can be written in normal form and falls within the general case
we considered in the previous chapter.
In particular, 
theorem~\ref{th:edo_lineare_ordine_n}
tells us that the solution space is an affine linear space of dimension $n$.

If $g=0$ we will say that the equation is linear homogeneous (still with
constant coefficients)
\index{equation!differential!linear!with constant coefficients}
\begin{equation}\label{eq:edo_lineare_omogenea_coeff_costanti}
    a_n \cdot u^{(n)}(x) + a_{n-1} \cdot u^{(n-1)}(x) + \dots + a_1 \cdot u'(x)
  + a_0 \cdot u(x) = 0
%   \sum_{k=0}^n a_k \cdot u^{(k)}(x) = 0.
\end{equation}
In this case, the set of solutions is a vector space of dimension $n$.
We observe that the global existence and uniqueness theorem guarantees that the
solutions
of the linear differential equation with constant coefficients are functions of 
class $C^n$ defined on the entire $\RR$. 
But the constant coefficients are of class $C^\infty$, thus the greater
regularity of the solutions
(theorem~\ref{th:maggiore_regolarita}) 
tells us, in this case, that the solutions must be functions of class
$C^\infty$.

For purely algebraic reasons, it will be useful to consider complex-valued
functions.
Recall that if $u(x)$ is a complex-valued function,
then it can be written as $u(x) = f(x) + i g(x)$ with $f$ and $g$ real-valued
functions.
Then we have $u'(x) = f'(x) + ig'(x)$.
We denote by $D \colon C^\infty(\RR, \CC)\to C^\infty(\RR, \CC)$ the derivative
operator: $D u = u'$.

We explicitly observe that even when $\lambda \in \CC$, it holds,
as in the case $\lambda\in \RR$
\begin{equation}\label{eq:499375}
  D e^{\lambda x} = \lambda e^{\lambda x}.
\end{equation}
One could perform a direct verification by reducing to the functions $\sin$ and
$\cos$
via Euler's formulas, but this
is actually a fundamental property of the complex exponential
that directly follows from theorem~\ref{th:exp_complesso} (we only need $h,x\in
\RR$, but in reality
this is true even if $h,x\in \CC$):
\[
 \lim_{h\to 0} \frac{e^{\lambda(x+h)}-e^{\lambda x}}{h}
 = \lim_{h\to 0} e^{\lambda x}\frac{e^{\lambda h}-1}{h}
 = \lambda e^{\lambda x} \lim_{h\to 0}\frac{e^{\lambda h}-1}{\lambda h}
 = \lambda e^{\lambda x}.
\]

If $P\in \RR[x]$ is a polynomial of degree $n$,
we can evaluate $P$ on the derivative operator
$D\colon C^\infty(\RR,\CC)\to C^\infty(\RR,\CC)$.
This makes sense because on the space $C^\infty(\RR,\CC)$
the operations of a real vector space
(addition and scalar multiplication) are defined, along with the operation
of multiplication $D\cdot D$, and thus the power
$D^k$ is understood as the composition of the linear operator
with itself.
In practice, the operator
$P(D)\colon C^\infty(\RR,\CC) \to C^\infty(\RR,\CC)$
is defined as follows
\[
  P(D)[u] = \sum_{k=0}^n a_k D^k u = \sum_{k=0}^n a_k u^{(k)}.
\]
Note that $D^0 = I$ is the identity, but keeping the convention
$D^0 = 1$ valid, we will occasionally write
$\lambda$ instead of $\lambda I$ when $\lambda\in \RR$.

The linear homogeneous equation with constant
coefficients~\eqref{eq:edo_lineare_omogenea_coeff_costanti}
can be more expressively written in the form
\[
  P(D)[u] = 0
\]
Since $a_n\neq 0$, we will have $\deg P = n$.

\subsection{Homogeneous equations}

\begin{theorem}
\mymark{***}
Let $P$ be a polynomial and let $\lambda\in \CC$ be a root of $P$ with
multiplicity $m$.
Then if $p(x)$ is a polynomial (with complex coefficients) of degree less than
$m$,
the function $u(x) = p(x) e^{\lambda x}$ is a (complex) solution of the
differential equation
\[
   P(D) [u] = 0.
\]
\end{theorem}
%
\begin{proof}
\mymark{***}
If the polynomial $P(t)$ has a root $\lambda$ with multiplicity $m$, it means
that $P(t)$ is divisible by $(t-\lambda)^m$, i.e., there exists a polynomial
$R(t)$ such that
\[
  P(t) = R(t)\cdot (t-\lambda)^m.
\]

But then we can also decompose the differential equation:
\[
 P(D) [u] = R(D)[(D - \lambda)^m u].
\]
If $u(x) = p(x) e^{\lambda x}$, we have,
using~\eqref{eq:499375},
\begin{align*}
  (D-\lambda) u(x) &= Du(x) -\lambda u(x) =
    p'(x) e^{\lambda x} + p(x) \lambda e^{\lambda x} - \lambda p(x) e^{\lambda
  x} \\
  &= p'(x) e^{\lambda x},
\end{align*}
from which
\[
  (D-\lambda)^m u(x) = p^{(m)}(x) \cdot e^{\lambda x}.
\]
Since the derivative of a polynomial is a polynomial of lower degree,
if the polynomial $p$ has degree less than $m$, it follows that $p^{(m)}=0$.
Thus, as we wanted to prove,
\[
 P(D) [u] = R(D) [(D-\lambda)^m u] = R(D) [0] = 0.
\]
\end{proof}

\begin{theorem}[independence of complex fundamental solutions]
\mymark{***}
The subset $\B$ of $C^\infty(\RR, \CC)$ given by the functions $u(x)$ of the
form
\[
   u(x) = x^m e^{\lambda x}
\]
as $m\in \NN$, and $\lambda \in \CC$ vary, is a linearly independent subset
of the vector space $C^\infty(\RR,\CC)$ over the field $\CC$.
\end{theorem}
%
\begin{proof}
It is a matter of showing that if a finite linear combination of such functions
is identically zero, then all the coefficients are zero.

Suppose, for contradiction, that there exist
$u_1, \dots, u_N\in \B$
\[
  u_k(x) = x^{m_k} e^{\lambda_k x}, \qquad k=1, \dots, N
\]
such that
\[
\sum_{k=1}^N c_k x^{m_k} e^{\lambda_k x} = 0
  \qquad \forall x \in \RR
\]
for some choice of coefficients $c_k \in \CC$ not all zero.
By summing together the monomials that multiply the exponentials with the same
coefficient, we can rewrite the previous relation in the form:
\begin{equation}\label{eq:4656978}
  \sum_{j=1}^M P_j(x) e^{\lambda_j x} = 0 \qquad \forall x\in \RR
\end{equation}
with the coefficients $\lambda_j$ all different from each other and with
$P_j$ non-zero polynomials.

We have already observed that $(D-\lambda)P(x)e^{\lambda x}= P'(x)e^{\lambda x}$
while if $\lambda\neq \mu$, we have $(D-\lambda)P(x) e^{\mu x} = Q(x) e^{\mu x}$
with $Q$ a polynomial of the same degree as $P$ because:
\[
(D-\lambda) P(x) e^{\mu x} = (P'(x) + (\mu-\lambda) P(x)) e^{\mu x}.  
\]
Let $d_j = \deg P_j$. We apply to equation~\eqref{eq:4656978}
the operator $(D-\lambda_1)^{d_1}$.
What is obtained is:
\begin{equation}\label{eq:4656979}
    k_1 e^{\lambda_1 x} + \sum_{j=2}^M Q_j(x) e^{\lambda_j x} = 0 \qquad \forall
  x\in \RR
\end{equation}
where $k_1$ is the $d_1$-th derivative of the polynomial $P_1$.
Since $P_1$ is a polynomial of degree $d_1$, its
$d_1$-th derivative is a polynomial of degree $0$ not zero,
thus it is a constant $k_1 \neq 0$.
The polynomials $Q_j$ have instead the same degree as the $P_j$ because we have
seen
that if $\lambda\neq \mu$, the degree of the polynomial does not change.

Now we apply in sequence to equation~\eqref{eq:4656979} the operators
$(D-\lambda_j)^{d_j+1}$ for $j=2,\dots, M$.
Applying these operators, the coefficient $k_1$ will change, but it will remain
non-zero (let's call it $k\neq 0$) while all the other polynomials
will be annulled because their respective polynomials are derived one more time
than
their degree. We will obtain:
\[
  k e^{\lambda_1 x} = 0 \qquad \forall x\in \RR.
\]
But this is absurd because for $x=0$, we obtain $k=0$, which we have excluded.
\end{proof}

The previous results allow us to determine all complex solutions of linear 
homogeneous equations with constant coefficients. Indeed, for each complex root
$\lambda$ of the associated polynomial
$P(t) = a_n t^n + \dots + a_1 t + a_0$, if $m$ is the multiplicity of $\lambda$,
the functions 
\[
  e^{\lambda x}, x e^{\lambda x}, \dots, x^{m-1} e^{\lambda x}  
\]
are $m$ independent solutions, and as $\lambda$ varies among the roots of $P$,
we obtain
a total of $n$ independent solutions. Since 
we know that the solution space has dimension $n$, the linear combinations of
these
solutions give us the set of all solutions.

If all the roots of the polynomial $P$ are real, the functions $e^{\lambda x}$
are indeed
real solutions, and, thanks to the following lemma, we can assert that all real
solutions
are obtained by making linear combinations with real coefficients of these
solutions.

\begin{lemma}\label{lm:combinazioni_reali}
Let $u_1(x), \dots, u_n(x)$ be independent real-valued functions
and let 
\[
  h(x) = c_1 \cdot u_1(x) + \dots + c_n \cdot u_n(x)
\]
be a linear combination with complex coefficients.
Then $h(x)$ is real-valued if and only if the coefficients $c_1, \dots, c_n$ 
are also real.
\end{lemma}
%
\begin{proof}
It is obvious that if the coefficients $c_1,\dots,c_n$ are real and the
functions $u_1, \dots, u_n$ are real,
making a linear combination of them results in a real function.

To prove the converse, it is necessary to know that the functions are also
independent. Indeed,
if $h(x)$ is real, we can write $\bar h(x) - h(x) = 0$, that is, since the
conjugate
passes through the sum and product and since $\overline u_k(x) = u_k(x)$, we
have
\[
 (\bar c_1 - c_1) u_1(x) + \dots + (\bar c_n-c_n) u_n(x)=0.
\]
But the zero linear combination can only be obtained with all coefficients zero 
(it is the definition of linear independence), thus $\bar c_k = c_k$ for each
$k$, meaning
the coefficients are real.
\end{proof}

\begin{example}
Solve the equation
\[
  u^{(5)} - 2 u''' + u' = 0.  
\]
\end{example}
%
\begin{proof}[Solution.]
The polynomial associated with the equation is $P(\lambda) = \lambda^5 -
2\lambda^3 + \lambda$
and can be factored as follows:
\[
  \lambda^5 - 2\lambda^3 + \lambda
  = \lambda(\lambda^4-2\lambda^2+1)
  = \lambda(\lambda^2-1)^2
  = \lambda(\lambda-1)^2(\lambda+1)^2.
\]
Thus, we have three distinct roots $\lambda_1=0$, $\lambda_2=1$, $\lambda_3=-1$
with multiplicities
$m_1=1$, $m_2=2$, $m_3=2$.
A basis of solutions is therefore given by the five independent functions
\[
 1, e^x, xe^x, e^{-x}, xe^{-x}.
\]
All solutions can thus be written in the form
\[
  u(x) = c_1 + (c_2 + c_3 x)e^x + (c_4 + c_5 x) e^{-x}  
\]
with $c_1, \dots, c_5$ arbitrary constants.
\end{proof}

However, if the roots of the polynomial are complex, we are still interested in
identifying
the real solutions.
To do this, we observe that if $\lambda = \alpha + i \beta $ is a complex root
of a polynomial $P$ with real coefficients, then $\bar \lambda =\alpha - i
\beta$ is also a complex root with the same multiplicity as $\lambda$
(see the proof of theorem~\ref{th:fattorizzazione_polinomio_reale}). 
But then we observe that the space generated by the linear combinations with
complex coefficients
of two solutions corresponding to the two conjugate coefficients:
\begin{align*}
  c_1 x^k e^{\lambda x} + c_2  x^k e^{\bar\lambda x}
  &= c_1 x^k e^{\alpha x+i \beta x} + c_2 x^k e^{\alpha x - i \beta x}\\
  &= x^k e^{\alpha x}\enclose{(c_1+c_2) \cos \beta x + i (c_1-c_2) \sin \beta x}
\end{align*}
is also a linear combination with complex coefficients of the two real
functions:
\begin{equation}\label{eq:489571798}
  x^k e^{\alpha x} \cos(\beta x), \qquad
  x^k e^{\alpha x} \sin(\beta x)
\end{equation}
Thus, by lemma~\ref{lm:combinazioni_reali}, 
the real solutions can be written as linear combinations with constant
coefficients
of these latter functions.

In practice, the method for determining the solutions is therefore very simple. 
Each real root $\lambda$ of the polynomial $P$ associated with the equation
gives rise
to the following independent solutions:
\[
  e^{\lambda x}, x e^{\lambda x}, \dots, x^{m-1} e^{\lambda x}
\]
where $m$ is the multiplicity of $\lambda$ as a root of the polynomial.
For each pair of complex roots $\lambda = \alpha \pm i \beta$ with multiplicity
$m$,
we will have the following independent solutions:
\begin{gather*}
    e^{\alpha x}\cos(\beta x), x e^{\alpha x}\cos(\beta x), \dots x^{m-1}
  e^{\alpha x}\cos(\beta x),\\
    e^{\alpha x}\sin(\beta x), x e^{\alpha x}\sin(\beta x), \dots x^{m-1}
  e^{\alpha x}\sin(\beta x).
\end{gather*}
In total, we always have $n=\deg P$ independent solutions.

\begin{example}
Determine all solutions of the differential equation
\[
  u^{(5)}(x) + 2 u'''(x) + u'(x) = 0.
\]
\end{example}
%
\begin{proof}[Solution]
The equation can be written in the form
\[
  P(D) [u] = 0
\]
with $P(\lambda) = \lambda^5 + 2 \lambda^3+\lambda$. 
We can factor the polynomial $P$ in the complex field:
\[
\lambda^5 + 2\lambda^3 +\lambda = \lambda(\lambda^2+1)^2 = \lambda (\lambda+i)^2(\lambda-i)^2.
\]
The polynomial has one root $\lambda_0=0$ with multiplicity one and two
conjugate complex roots
$\lambda_1 = i$, $\lambda_2=-i$ with multiplicity two.
It follows that the following functions must be independent real solutions of
the equation:
\begin{equation}\label{eq:4995566}
 1, \qquad
 \cos x, \qquad
 x \cos x, \qquad
 \sin x, \qquad
 x \sin x
\end{equation}
Every real solution can thus be written in the form:
\begin{equation}\label{eq:499355456}
  u(x) = c_1 + c_2 \cos x + c_3 \sin x + c_4 x \cos x + c_5 x \sin x
\end{equation}
with $c_1, c_2, \dots, c_5 \in \CC$ appropriate real coefficients.
\end{proof}

% Formalizing in general
% the method used in the previous exercise, we can
% give the following statement.
% 
% \begin{theorem}[solutions of the linear homogeneous equation with constant coefficients]
% \mymark{**}%
% If $P(t)$ is a polynomial with real coefficients.
% Every solution
% $u\colon \RR\to \RR$ of the differential equation
% \[
%   P(D) [u] = 0
% \]
% can be written in the form
% \begin{equation}
% \label{eq:4725549774}
%   u(x) = \sum_{k=1}^N \sum_{j=1}^{n_k} c_{kj} x^j e^{\lambda_k x}
%         + \sum_{k=1}^M \sum_{j=1}^{m_l} x^j e^{\alpha_k x} (a_{kj} \cos(\beta_k x) + b_{kj}\sin(\beta_k x))
% \end{equation}
% where $\lambda_1, \dots, \lambda_N$ are the real roots of the polynomial $P(t)$, $n_1, \dots, n_N$ are the respective multiplicities,
% $\alpha_k \pm i \beta_k$ are the conjugate complex (non-real) roots of the polynomial $P$ with $k=1,\dots, M$ each with multiplicity $m_k$ and finally $c_{kj}, a_{kj}, b_{kj}\in \RR$ are arbitrary constants.
% \end{theorem}
% %
% \begin{proof}
% We know that the complex functions
% \begin{equation}\label{eq:739647}
%   x^j e^{\lambda_k x}, \qquad
%   x^j e^{(\alpha_k+i\beta_k)x}
% \end{equation}
% are solutions of the equation if $j$ is less than the multiplicity of the corresponding real root $\lambda_k$ or complex $\alpha_k + i \beta_k$.
% Thus
% \[
%   u_{k,j}(x) = x^j e^{\lambda_k x}
% \]
% is a solution if $j<n_k$.
% For the complex roots, we apply Euler's formula:
% \begin{align}\label{eq:49544467}
% x^j e^{\alpha_k x}\cos (\beta_k x) &=
%   \frac 1 2 x^j e^{(\alpha_k +i \beta_k)x}
%   + \frac 1 2 x^j e^{(\alpha_k - i \beta_k)x}\\
% x^j e^{\alpha_k x}\sin (\beta_k x) &=
%     \frac 1 {2i} x^j e^{(\alpha_k +i \beta_k)x}
%     - \frac 1 {2i} x^j e^{(\alpha_k - i \beta_k)x}
% \end{align}
% thus even the functions
% \[
%   v_{kj}(x) = x^j e^{\alpha_k x}\cos x, \qquad
%   w_{kj}(x) = x^j e^{\alpha_k x}\sin x
% \]
% being a linear combination of solutions, are also
% solutions.
% These solutions are also independent because
% the functions in \eqref{eq:739647} are (by the previous theorem) and can be reduced to \eqref{eq:49544467}
% via Euler's formula.
% 
% We have thus found a set formed by
% \[
%   n = \sum_{k=1}^N n_k + 2 \sum_{k=1}^M m_k
% \]
% independent real solutions. The number $n$ is equal to the sum of the multiplicities of the 
% (real and complex) roots of the polynomial $P$ and thus, by the fundamental theorem of algebra, 
% coincides with the degree of $P$. We then know that the solution space of $P(D)[u]=0$ has 
% dimension $n$ and therefore we have found a basis of all real solutions. 
% This means that every solution can be written in the form~\eqref{eq:4725549774}.
% \end{proof}
% 
\begin{remark}[phase of oscillating solutions]
When writing the linear combinations of the functions 
\eqref{eq:489571798},
it may be useful to observe
that every linear combination of sines and cosines
with the same frequency $\beta$ is a single wave
possibly phase-shifted by an angle $\theta$:
\[
  a \cos (\beta x) + b \sin (\beta x)
  =
  \rho \cos(\beta x - \theta)
\]
where $\rho = \abs{a+ib} = \sqrt{a^2+b^2}$ and 
$\theta=\arg(a+ib)$. 
Indeed,
setting $a=\rho\cos \theta$, $b=\rho \sin \theta$, we have,
thanks to the addition formula:
\[
   a \cos (\beta x) + b \sin (\beta x)
   = \rho\enclose{\cos \theta \cos(\beta x) + \sin \theta \sin (\beta x)}
   = \rho \cos(\beta x - \theta).
\]
\end{remark}

\subsection{Non-homogeneous linear equations}

Theorem~\ref{th:edo_lineare_ordine_n}
tells us that the solutions of a 
non-homogeneous linear equation (whether with constant coefficients or not)
\[
  L[u]=g
\]
can all be written in the form $u = u_* + \ker L$ where 
$u_*$ is any fixed solution of the non-homogeneous equation 
and $\ker L$ is the set of solutions of the associated homogeneous equation 
\[
  L[u] =0. 
\]
Then, assuming we know how to solve the homogeneous equation 
(which we do if the coefficients are constant),
to find all the solutions of the non-homogeneous equation, it will suffice 
to find just one particular solution $u_*$.

In the following, we will see two methods for determining a particular solution 
when we have a non-homogeneous equation with constant coefficients. 
The first method (of similarity) applies when the known term $g(x)$ takes a
specific form,
the second method (of variation of constants) applies in general, but it might
be
more laborious than the first.

\begin{theorem}[method of similarity]
\mymark{***}
Consider the complex non-homogeneous equation
\begin{equation}\label{eq:3976734956}
 P(D) [u] = q(x) e^{\lambda x}.
\end{equation}
with $P$ and $q$ polynomials with complex coefficients.
Let $m$ be the multiplicity of $\lambda$ as a root of $P$ ($m=0$ if $\lambda$ is
not a root of $P$).
Then a particular solution of~\eqref{eq:3976734956} can be written in the form
\begin{equation}
\label{eq:945396}
u_*(x) = p(x) x^m e^{\lambda x}
\end{equation}
with $p$ an (unknown) polynomial of degree equal to the degree of $q$.
\end{theorem}
%
\begin{proof}
If $u_*(x) = x^m p(x) e^{\lambda x}$, we want to understand how the operator
$P(D)$ acts on $u_*$, in particular
given any polynomial $q$, we want to understand
if there exists a polynomial $p$, of the same degree as $q$, such that applying
$P(D)$ to $u_*(x)$ results in $q(x) e^{\lambda x}$.

Factor the polynomial $P(z)$:
\[
  P(z) = (z-\lambda_1)^{m_1}\dots (z-\lambda_n)^{m_n}
\]
where $\lambda_1, \dots, \lambda_n\in \CC$ are the distinct roots of $P(z)$ and
$m_1,\dots,m_n\in \NN$ are the respective multiplicities. Thus, the operator
$P(D)$ can be factored in the corresponding form:
\[
   P(D) = (D-\lambda_1)^{m_1} \dots (D-\lambda_n)^{m_n}.
\]

Let $V^n_m$ be the vector space of polynomials of the form $x^m p(x)$ with $\deg
p<n$.
A basis of $V^n_m$ are the monomials $x^m, x^{m+1}, \dots, x^{m+n-1}$, and thus
$V^n_m$ is a
vector space of dimension $n$.

The proof of the theorem is based on the observation of how the operator
$(D-\lambda_k)^{m_k}$
acts on functions of the type $u(x) = p(x) e^{\lambda x}$ with $p\in V^n_m$. In
general, if $p$ is a
polynomial, we have
\[
 (D-\lambda_k) [p(x)\cdot e^{\lambda x}]
 = \Enclose{(\lambda -\lambda_k)p(x) + p'(x)} \cdot e^{\lambda x}.
\]
We must then distinguish two cases. If $\lambda\neq \lambda_k$, then the
function $p(x) e^{\lambda x}$ is transformed into the function $q(x) e^{\lambda
x}$ where $q$ is a polynomial of the same degree as $p$, since the highest
degree term of $p$ is multiplied by $\lambda-\lambda_k$ while the polynomial
$p'$ has a degree lower than that of $p$ and therefore does not influence the
highest degree term. If we call $T_k\colon V^n_0 \to V^n_0$ the linear operator
that maps the polynomial $p(x)$ to the polynomial $q(x)=(\lambda-\lambda_k) p(x)
+ p'(x)$, we observe that $T_k$ is injective because, since $T_k$ preserves the
degree of the polynomial, the only polynomial that can go to zero is the zero
polynomial. Thus, by the rank theorem, $T_k$ is bijective, and consequently,
$T_k^{m_k} \colon V^n_0 \to V^n_0$ is bijective.

If instead $\lambda = \lambda_k$, we now observe that the operator $T_k$ 
is nothing but the derivative: $q(x) = p'(x)$. 
And if $m=m_k$ is the multiplicity of $\lambda$ as a root of $P$, 
the operator $T_k^m$ is nothing but the $m$-th derivative. 
This operator is not invertible on $V^n_0$ because it sends to zero all 
polynomials of degree less than $m$. But if we start from a polynomial in
$V^n_m$,
then the operator $T_k^m\colon V^n_m \to V^n_0$ becomes invertible: 
indeed, in $V^n_m$, there are polynomials of the form 
$a_m x^m + a_{m+1}x^{m+1} + \dots + a_{n+m}x^{n+m}$
and taking the $m$-th derivative can only result in the zero polynomial if all 
the coefficients are zero.
It follows that, if $\lambda = \lambda_k$, 
the operator $T_k^m\colon V^n_m \to V^n_0$ is invertible.

We have reached the conclusion. If we take a polynomial $p(x)$ in $V^n_m$ and
apply $P(D)$
to the corresponding function $u(x) =p(x) e^{\lambda x}$,
we can sequentially apply the operators $(D-\lambda_k)^{m_k}$. 
These act on the function $u$ by modifying the polynomial $p$.
If $\lambda$ is a root of $P$, at the first step, we apply the operator
$(D-\lambda_k)^{m_k}$
with $\lambda_k=\lambda$ and $m_k=m$ so that the polynomial $p\in V^n_m$ is
transformed,
in a one-to-one manner, into a polynomial of $V^n_0$. 
Then, I apply all the other operators $(D-\lambda_k)^{m_k}$ with $\lambda_k \neq
\lambda$
transforming the polynomial of $V^n_0$ into another polynomial of $V^n_0$, but
always in a one-to-one manner.
It follows that in the end, we obtain a one-to-one map between $V^n_m \to
V^n_0$, and whatever
$q\in V^n_0$, we know there exists a $u\in V^n_m$ that is mapped to $q$.
\end{proof}

\begin{remark}
Note that in the previous proof, we investigated the
structure of the differential operator $D$ with constant
coefficients.
What we did is actually the process of
Jordan decomposition of a linear operator.
The operator $D$
has as eigenvectors the functions $e^{\lambda_k x}$ with $\lambda_k$
roots of the polynomial $P$.
It is observed, however, that functions of the form
$x^j e^{\lambda_k x}$ are \emph{generalized eigenvectors}
since iterating the operator $D-\lambda_k$ sends them to an eigenvector
and they are thus a basis of the eigenspace relative to the eigenvector
$\lambda_k$.

The basis formed by the functions of the form $\frac{x^j}{j!} e^{\lambda_j x}$
is
a \emph{fan basis}, and the matrix representing the operator
$D$ turns out to be a block matrix in Jordan form: the
eigenvalues are on the diagonal, and a row of $1$s is found above the
diagonal.

The kernel of the operator $P(D)$ decomposes as a direct sum
of the kernels of the operators $(D-\lambda_k)^{m_k}$.
\end{remark}

\begin{definition}[Wronskian]
Let $u_1, \dots, u_n$ be functions of class $C^{n-1}$.
The matrix
\begin{equation}\label{eq:wronksiano}
  W(x) =
    \begin{pmatrix}
    u_1(x) & u_2(x) & \dots & u_n(x) \\
    u_1'(x) & u_2'(x) & \dots & u_n'(x) \\
    \vdots & & & \vdots\\
    u_1^{(n-1)}(x) & u_2^{(n-1)}(x) & \dots & u_n^{(n-1)}(x)
    \end{pmatrix}
\end{equation}
is called the \emph{Wronskian matrix}.
\mymargin{Wronskian matrix}%
\index{Wronskian matrix}%
\index{matrix!Wronskian}%
\index{Wronskian}%
The determinant of this matrix
$w(x) = \det W(x)$ is called the
\emph{Wronskian determinant}%
\mymargin{Wronskian determinant}\index{determinant!Wronskian}.
\end{definition}

\begin{theorem}%
\label{th:wronksiano}%
Let $I\subset \RR$ be a non-empty open interval and let $a_j\in C^0(I,\CC)$.
Let $u_1, \dots, u_n$ be functions of class $C^n$ solutions
of the homogeneous linear equation in normal form:
\begin{equation}\label{eq:092784}
  u^{(n)} = \sum_{j=1}^n a_j(x) u^{(j-1)}(x).
\end{equation}
Let $W(x)$ be the Wronskian matrix~\eqref{eq:wronksiano}.
Then the following properties are equivalent:
\begin{enumerate}
\item \label{item:4736201} the functions $u_1,\dots, u_n$ are linearly independent;
\item \label{item:4736202} for every $x\in I$, we have $\det W(x)\neq 0$;
\item \label{item:4736203} there exists $x\in I$ such that $\det W(x)\neq 0$.
\end{enumerate}
\end{theorem}

\begin{proof}
Obviously, \ref{item:4736202} implies \ref{item:4736203}.

We prove that \ref{item:4736203} implies \ref{item:4736201}.
Passing to the contrapositive implication, we must show that
if $u_1,\dots, u_n$ are linearly dependent, then for every $x\in I$, we have
$\det W(x)=0$. But if the functions are dependent, it means there exists $\vec
\lambda \neq 0$ such that
\begin{equation}\label{eq:3675893}
  \lambda_1 u_1(x) + \dots + \lambda_n u_n(x) = 0 \qquad \forall x\in I.
\end{equation}
Differentiating, we obtain:
\[
\begin{cases}
  \lambda_1 u_1(x) + \dots + \lambda_n u_n(x) = 0\\
  \lambda_1 u'_1(x) + \dots + \lambda_n u'_n(x) = 0\\
  \dots\\
  \lambda_1 u^{(n-1)}_1(x) + \dots + \lambda_n u^{(n-1)}_n(x) = 0
\end{cases}
\]
or
\[
  W(x) \, \vec \lambda = 0 \qquad \forall x\in I.
\]
But then the matrix $W(x)$ is not invertible, and therefore $\det W(x)=0$, as we
wanted to prove.

Finally, we prove that \ref{item:4736201} implies \ref{item:4736202}
that is, if $v_1, \dots, v_n$ are independent, then $\det W(x) \neq 0$ for every
$x\in I$.
Fixed $x\in I$, consider the functional $J(v) = (v(x),v'(x), \dots,
v^{(n-1)}(x))$ so that we have
\[
  W(x) = \enclose{J(u_1), \dots, J(u_n)}.
\]
According to theorem~\ref{th:edo_lineare_ordine_n}, we recall that $J$ turns out
to be an isomorphism of vector spaces, so if $v_1,\dots,v_n$ are independent,
then $J(v_1),\dots, J(v_n)$ must be independent.
But since the columns of the matrix $W(x)$ are independent, it means that $\det
W(x)\neq 0$, as we wanted to prove.
\end{proof}

\begin{theorem}[method of variation of arbitrary constants]
\mymark{***}%
Consider a generic non-homogeneous linear differential equation of order $n$ in
normal form:
\begin{equation}\label{eq:395467435}
  u^{(n)}(x) = \sum_{k=0}^{n-1} a_k(x) u^{(k)}(x) + b(x)
\end{equation}
where $a_0, \dots, a_{n-1}, b\in C^0(A,\CC)$ are functions defined on an open
set $A\subset \RR$
and we seek a solution $u\in C^n(A,\CC)$.

If $v_1, \dots, v_n\in C^n(A,\CC)$ are independent solutions of the associated
homogeneous equation:
\begin{equation*}
  u^{(n)}(x) = \sum_{k=0}^{n-1} a_k(x) u^{(k)}(x)
\end{equation*}
then
there exist functions $c_k\in C^1(A,\CC)$
for $k=0,1, \dots, n-1$
such that
\begin{equation}\label{eq:02154676}
  \begin{cases}
    c'_1 v_1 + \dots + c'_n v_n = 0 \\
    c'_1 v_1' + \dots + c'_n v_n' = 0 \\
    \vdots\\
    c'_1 v_1^{(n-2)} + \dots + c'_n v_n^{(n-2)} = 0 \\
    c'_1 v_1^{(n-1)} + \dots + c'_n v_n^{(n-1)} = b(x).
  \end{cases}
\end{equation}
And in such a case, the function
\[
 u(x) = \sum_{k=1}^n c_k(x) v_k(x)
\]
solves the non-homogeneous equation~\eqref{eq:395467435}.
\end{theorem}
%
\begin{proof}
First, observe that if there exist functions $c_k$ that
satisfy the system~\eqref{eq:02154676}, then the function $u = c_1
\cdot v_1 + \dots + c_n \cdot v_n$ solves the non-homogeneous
equation. Indeed, from the product differentiation formula, we have:
\[
  u' = \sum_{k=1}^n c_k v_k' + \sum_{k=1}^n c_k' v_k
     = \sum_{k=1}^n c_k v_k'
\]
since the second term is zero by hypothesis~\eqref{eq:02154676}.
But then we can proceed with the derivatives:
\[
  u'' = \sum_{k=1}^n c_k v_k'' + \sum_{k=1}^n c_k' v_k'
      = \sum_{k=1}^n c_k v_k''
\]
up to the $(n-1)$-th derivative:
\[
  u^{(n-1)} = \sum_{k=1}^n c_k v_k^{(n-1)}.
\]
For the last derivative, we have:
\[
  u^{(n)} = \sum_{k=1}^n c_k v_k^{(n)} + \sum_{k=1}^n c_k' v_k^{(n-1)}
          = \sum_{k=1}^n c_k v_k^{(n)} + b.
\]
But then it is observed that we have
\begin{align*}
 u^{(n)} + \sum_{j=0}^{n-1} a_j u^{(j)}
  &= \sum_{k=1}^n c_k v_k^{(n)} + b + \sum_{j=0}^{n-1} a_j \sum_{k=1}^n c_k
 v_k^{(j)} \\
  &= \sum_{k=1}^n c_k \Enclose{v_k^{(n)} + \sum_{j=0}^{n-1} a_j v_k^{(j)}} + b =
 b
\end{align*}
 since the functions $v_j$, by hypothesis, solve the homogeneous equation.

 It remains only to verify that there exist functions $c_k\in
 C^1(A,\CC)$ that satisfy the system~\eqref{eq:02154676}.
 It is observed that the system~\eqref{eq:02154676} is written in the form
 \[
   W(x) \vec c'(x) = \vec b(x)
 \]
 where $W(x)$ is the Wronskian matrix $W(x)$ associated
 with the solutions $v_1, \dots, v_n$, $\vec c(x) = (c_1(x), \dots,
 c_n(x))$, and $\vec b(x) = (0, \dots, 0, b(x))$.
 Since $v_1, \dots, v_n$ are independent by hypothesis,
 by theorem~\ref{th:wronksiano},
 we know that
 $\det W(x)\neq 0$ for every $x$, and therefore the system $W(x) \vec d(x) =
 \vec b(x)$
 admits a unique solution $d(x)$ for every $x\in I$.
 Since $W(x)$ and $\vec b(x)$ have continuous coefficients,
  the solution $\vec d(x) = W(x)^{-1} \vec b(x)$ must also be continuous (that
 the coefficients of the matrix $W(x)^{-1}$ are continuous can be seen, for
 example, by writing $W(x)^{-1}$ using the matrix of cofactors and exploiting
 the fact that the determinant is a continuous function)
 and therefore, fixed a point $x_0\in I$,
 we can define
 \[
   c_k(x) = \int_{x_0}^x d_k(t)\, dt
 \]
 to determine the functions $c_k$.
\end{proof}

\begin{exercise}
Find all solutions of the non-homogeneous linear equation with
constant coefficients:
\[
  u''(x) + u(x) = \frac{1}{\cos x}.    
\]
\end{exercise}
\begin{proof}[Solution.]
The associated polynomial is $P(\lambda) = \lambda^2 + 1$ which 
has two conjugate complex roots $\pm i$.
Two independent solutions of the homogeneous equation $u''+u=0$ are thus 
(as known):
\[
 u_1(x) = \cos x, \qquad u_2(x) = \sin x.
\]
The general solution of the homogeneous equation is thus of the form:
\[
  u(x) = A \cos x + B \sin x.  
\]

To find a particular solution of the non-homogeneous equation, 
we use the method of variation of constants, i.e.,
we seek a solution of the form:
\[
  u_*(x) = A(x) \cos x + B(x) \sin x.
\]
The coefficients $A(x)$ and $B(x)$ are determined by solving the system 
\[
\begin{cases}
  A'(x) \cdot \cos x &+ B'(x)\cdot \sin x = 0, \\
  A'(x) \cdot (-\sin x) &+ B'(x)\cdot \cos x = \frac{1}{\cos x}.
\end{cases}  
\]
We find therefore 
\[
 \begin{cases}
  A'(x) = -\frac{\sin x}{\cos x} \cdot B'(x) \\
  B'(x) = 1
 \end{cases}  
\]
from which, integrating, we find
\[
  A(x) = \ln\abs{\cos x}, \qquad B(x) = x
\]
and therefore 
\[
 u_*(x) = \cos x\cdot \ln \abs{\cos x } + x\sin x.  
\]

All solutions of the given equation are thus written as the sum of the 
particular solution of the non-homogeneous equation plus the general solution of
the
homogeneous equation:
\[
 u(x) = a \cos x + b \sin x + \cos x \cdot \ln \abs{\cos x} + x \sin x.  
\]
\end{proof}